\documentclass[10pt]{article}
\usepackage[T1]{fontenc}
\usepackage[utf8]{inputenc}
\usepackage[french]{babel}
\usepackage{amsmath,amssymb,mathrsfs,mathpazo}
\usepackage{geometry}
\geometry{a4paper,margin=2cm}
\usepackage{graphicx,array,tabularx,tabularray,colortbl,multirow,multicol,booktabs,enumitem,pifont,ulem,fancybox,diagbox,cancel}
\usepackage{pstricks,pstricks-add,pst-plot,pst-tree,pst-node,pst-text,pst-func,pst-math,pst-bspline,pst-3dplot,pst-eucl,pst-solides3d}
\let\pscalculate\undefined
\usepackage{pst-bezier}
\usepackage{tikz,pgfplots,tkz-tab}
\usetikzlibrary{arrows,arrows.meta,patterns,calc,positioning,trees,decorations.markings,shapes.geometric,shapes.misc,shapes.arrows,intersections,decorations.pathmorphing,decorations.pathreplacing,backgrounds,shadows,babel}
\pgfplotsset{compat=1.18}
\usepgfplotslibrary{fillbetween}
\usepackage[np]{numprint}
\usepackage{listings,makecell,pas-tableur,twoopt}
\let\barre\undefined
\usepackage[pstricks]{tablvar}
\definecolor{darkgreen}{RGB}{0,100,0}
\providecommand{\red}{\color{red}}
\providecommand{\blue}{\color{blue}}
\providecommand{\green}{\color{green}}
\providecommand{\black}{\color{black}}
\providecommand{\white}{\color{white}}
\providecommand{\gray}{\color{gray}}
\providecommand{\cyan}{\color{cyan}}
\providecommand{\magenta}{\color{magenta}}
\providecommand{\yellow}{\color{yellow}}
\providecommand{\violet}{\color{violet}}
\providecommand{\orange}{\color{orange}}
\providecommand{\euro}{\texteuro}
\providecommand{\EUR}{\texteuro}
\providecommand{\ieme}{\textsuperscript{e}}
\providecommand{\degres}{\ensuremath{^\circ}}
\providecommand{\degree}{\ensuremath{^\circ}}
\providecommand{\vv}[1]{\overrightarrow{#1}}
\providecommand{\Coord}[1]{#1}
\providecommand{\decosix}{}
\providecommand{\starredbullet}{\ensuremath{\bullet}}
\pagestyle{empty}
\setlength{\parindent}{0pt}
\usepackage{adjustbox}\DeclareUnicodeCharacter{2236}{:}\DeclareUnicodeCharacter{2190}{\ensuremath{\leftarrow}}
\providecommand{\R}{\mathbb{R}}
\providecommand{\N}{\mathbb{N}}
\providecommand{\Z}{\mathbb{Z}}
\providecommand{\Q}{\mathbb{Q}}
\providecommand{\D}{\mathbb{D}}
\providecommand{\C}{\mathbb{C}}
\providecommand{\e}{{\mathrm{e}}}
\providecommand{\vect}[1]{\overrightarrow{#1}}
\providecommand{\vectt}[1]{\overrightarrow{\mathrm{#1}}}
\providecommand{\vvt}[1]{\overrightarrow{\mathrm{#1}}}
\providecommand{\np}[2][]{{#2~#1}}
\providecommand{\vv}[1]{{\overrightarrow{#1}}}
\renewcommand{\d}{{\mathrm{d}}}
\providecommand{\degres}{{^\circ}}
\providecommand{\degre}{{^\circ}}
\providecommand{\blue}{\color{blue}}
\providecommand{\red}{\color{red}}
\providecommand{\enskip}{\quad}
\providecommand{\llbracket}{\lbrack\!\lbrack}
\providecommand{\rrbracket}{\rbrack\!\rbrack}
\providecommand{\euro}{{\text{€}}}
\providecommand{\no}{{\text{n°}}}
\providecommand{\interval}[2]{{[#1;#2]}}
\providecommand{\intervFO}[2]{{[#1;#2[}}
\providecommand{\dotfill}{\ldots}

\providecommand{\ds}{\displaystyle}
\providecommand{\textding}[1]{#1}
\providecommand{\Oij}{\(\mathrm{O};\vec{\imath},\vec{\jmath}\)}
\providecommand{\Oijk}{\(\mathrm{O};\vec{\imath},\vec{\jmath},\vec{k}\)}
\providecommand{\Ouv}{\(\mathrm{O};\vec{u},\vec{v}\)}
\providecommand{\intervFF}[2]{\ensuremath{[#1;#2]}}
\renewcommand{\euro}{\ifmmode\text{\texteuro}\else\texteuro\fi}
\DeclareUnicodeCharacter{00B0}{\ifmmode^{\circ}\else\textdegree\fi}
\DeclareUnicodeCharacter{20AC}{\ifmmode\text{\texteuro}\else\texteuro\fi}
% Énoncés : Éducation nationale. Transcription et corrigés : APMEP, Denis Vergès et contributeurs. Découpage et classement : CoopMaths, Sébastien Lozano, Éric Elter et contributeurs. Adaptation de la présentation : Maths Explorer.
% CC BY-SA : https://coopmaths.fr/www/about/
% https://coopmaths.fr/alea/static/bac/2023/tex/bac_2023_08_sujet1_caledonie_1_cor.tex
\begin{document}
\noindent\textit{APMEP / CoopMaths. CC BY-SA. Présentation adaptée par Maths Explorer.}\par\medskip
\medskip

\textbf{Partie A}

\medskip

\begin{enumerate}
	\item\
	\begin{center}
	\pstree[treemode=R,nodesepA=0pt,nodesepB=2.5pt,treesep = 1cm,levelsep=2.5cm]{\TR{}}
	{\pstree{\TR{$V$~~}\naput{\np{0.6}}}
		{\TR{$L$}\naput{\np{0.2}}
		\TR{$\overline{L}$}\nbput{\np{0.8}}
	}
	\pstree{\TR{$\overline{V}$~~}\nbput{\np{0.4}}}
		{\TR{$L$}\taput{$x$}
		\TR{$\overline{L}$}\nbput{$1-x$}
	}}
	\end{center}
    \item La probabilité qu'un client choisisse un bateau à voile et qu'il ne prenne pas l'option PILOTE est $P\left(V \cap \overline{L}\right)$.
	
	\[P\left(V \cap \overline{L}\right) = P(V) \times P_{V}\left(\overline{L}\right) = \np{0.6} \times \np{0.8} = \boxed{\np{0.48}}\]
	\item Les évènements $V$ et $\overline{V}$ forment une partition de l'univers, donc, d'après la formule des probabilités totales, on a :
	\[P(L) = P\left(V \cap L\right) + P\left(\overline{V} \cap L\right)\]
	On a $P\left(V \cap L\right) = \np{0.6} \times \np{0.2} = \np{0.12}$.
	
	Posons $P\left(\overline{V}\cap L\right) = p$. On a :
	\[P(L) = \np{0.12} + p \iff \np{0.42} = \np{0.12} + p \iff p = \np{0.42}-\np{0.12} = \boxed{\np{0.3}}\]
	On a donc bien $P\left(\overline{V} \cap L\right) = \np{0.3}$.
	\item D'après la formule des probabilités conditionnelles :
	\[P_{\overline{V}}(L) = \dfrac{P\left(\overline{V} \cap L\right)}{P\left(\overline{V}\right)} = \dfrac{\np{0.3}}{\np{0.4}} = \boxed{\dfrac{3}{4} = \np{0.75}}\]
	\item La probabilité qu'un client ait choisi un voilier, sachant qu'il a pris l'option PILOTE est, d'après la formule des probabilités conditionnelles :
	\[P_{L}(V) = \dfrac{P\left(V \cap L\right)}{P(L)} = \dfrac{\np{0.12}}{\np{0.42}} = \boxed{\dfrac{2}{7} \approx \np{0.29}}\]
\end{enumerate}

\medskip

\textbf{Partie B}

\medskip

\begin{enumerate}
	\item \[P\left(L \cap A\right) = P(L) \times P_{L}(A) = \np{0.42} \times \np{0.005} = \boxed{\np{0.0021}}\]
	\smallskip
	
	\[P\left(\overline{L} \cap A\right) = P\left(\overline{L}\right) \times P_{\overline{V}}(A) = (1-\np{0.42}) \times \np{0.12} = \boxed{\np{0.0696}}\]
	\item Les évènements $L$ et $\overline{L}$ forment une partition de l'univers, donc, d'après la formule des probabilités totales, on a :
	\[P(A) = P\left(L \cap A\right) + P\left(\overline{L} \cap A\right) = \np{0.0021} + \np{0.0696} = \np{0.0717}\]
La probabilité que le bateau loué par un client choisi au hasard subisse une avarie est donc \np{0.0717}.
	
Puisque l'entreprise loue \np{1000} bateaux, elle peut s'attendre à 72 avaries.
\end{enumerate}

\medskip

\textbf{Partie C}

\medskip

\begin{enumerate}
	\item Les paramètres de la loi binomiale suivie par $X$ sont $\boxed{n=40}$ et $\boxed{p=\np{0.42}}$.
	\item À l'aide de la calculatrice, on calcule $P(X \geqslant 15)$ :
	\[\boxed{P(X \geqslant 15) \approx \np{0.768}}\]
\end{enumerate}

\bigskip

\end{document}
